% !TeX program = lualatex-dev
% ===================================================================
% mwe_acessivel_book.tex — teste do livros_crp_acessivel_book.sty
%
% IMPORTANTE: compilar com lualatex-dev, não lualatex.
% A versão -dev do LaTeX usa os módulos latex-lab-dev, que contêm
% correções de tagging não disponíveis na versão estável do
% latex-lab. O ganho principal é a estruturação correta do sumário
% com <TOCI>/<Reference>/<Link>.
%
% Reproduz o MWE de scrbook anterior com a variante book, para
% comparação direta de árvore de tags e layout.
%
% Compilar com: lualatex-dev mwe_acessivel_book.tex (duas vezes)
% ===================================================================

\DocumentMetadata{
    lang        = pt-BR,
    pdfstandard = ua-2,
    pdfversion  = 2.0,
    testphase   = {phase-III,table,firstaid},
    tagging     = on,
}

\documentclass[11pt,a5paper,twoside,openright]{book}

\usepackage{livros_crp_acessivel_book}

\hypersetup{
    pdftitle  = {MWE — Teste de acessibilidade (book)},
    pdfauthor = {CRP-SP},
    pdflang   = {pt-BR},
}

\title{MWE — Teste de acessibilidade}
\author{CRP-SP}

\begin{document}

\frontmatter
\pagestyle{pretext}

% --- Folha de rosto ---
% TESTE DIAGNÓSTICO: usando \maketitle em vez de ambiente titlepage.
% O \title/\author já foram definidos no preâmbulo. \maketitle é o
% caminho canônico do book e o latex-lab-testphase-title foi escrito
% para suportá-lo explicitamente, enquanto titlepage livre produz
% <Part> genéricos e, pela hipótese atual, também descasa o
% contador begin/end de text-unit do tagpdf.
\maketitle

% Sumário via wrapper acessível (ver comentário em
% livros_crp_acessivel_book.sty, seção Sumário, para o motivo).
\TOCAcessivel

\mainmatter
\pagestyle{crp}

% ===================================================================
% PARTE I
% ===================================================================
\part{Primeira parte do teste}

% ===================================================================
% Capítulo 1
% ===================================================================
\chapter{Texto corrido e tipografia}

Este capítulo existe para verificar se a fonte Lora está sendo
carregada corretamente como fonte principal, se o espaçamento
entre linhas corresponde ao \texttt{\textbackslash onehalfspacing},
e se a hifenização portuguesa está ativa por meio do pacote
\texttt{babel}.

O próximo parágrafo contém um texto longo o suficiente para que
o algoritmo de quebra de linhas do \LaTeX{} precise hifenizar
pelo menos uma palavra composta, como
\emph{extraordinariamente} ou \emph{responsabilidade}, e assim
demonstrar que a hifenização está se comportando conforme o
esperado no contexto do idioma selecionado.

Aqui testamos \textbf{negrito}, \emph{itálico},
\textsf{fonte sem serifa (NEWJUNE)}, e o peso customizado
{\Light peso light da NEWJUNE} seguido por
{\media peso medium da NEWJUNE} e por fim
{\fontebook peso book da NEWJUNE}.

\section{Seção de primeiro nível}

Este texto está dentro de uma seção. Em PDF tagueado, deve
aparecer como \texttt{<H2>} (ou \texttt{<H>} com nível 2) na
árvore de estrutura.

\subsection{Subseção}

Texto da subseção. Deve virar \texttt{<H3>}.

\subsubsection{Sub-subseção}

Texto da sub-subseção.

\section{Outra seção}

Esta seção serve para verificar se há numeração visível (não
deve haver — \texttt{secnumdepth=-1}).

% ===================================================================
% Capítulo 2
% ===================================================================
\chapter{Elementos variados}

\section{Listas}

Lista com marcadores:
\begin{itemize}
    \item Primeiro item da lista.
    \item Segundo item, um pouco mais longo para testar quebra
          de linha dentro de item.
    \item Terceiro item.
\end{itemize}

Lista numerada:
\begin{enumerate}
    \item Primeiro.
    \item Segundo.
    \item Terceiro.
\end{enumerate}

\section{Tabela simples}

\begin{center}
\begin{tabular}{lcr}
    \toprule
    \textbf{Item} & \textbf{Quantidade} & \textbf{Valor} \\
    \midrule
    Livro         & 10                  & R\$ 250,00     \\
    Caderno       &  5                  & R\$  75,00     \\
    Caneta        & 20                  & R\$  40,00     \\
    \bottomrule
\end{tabular}
\end{center}

\section{Link e URL}

\url{https://www.crpsp.org.br}

E um link com texto customizado:
\href{https://www.crpsp.org.br}{site do CRP-SP}.

% ===================================================================
% Capítulo 3
% ===================================================================
\chapter{Página final}

Último capítulo do MWE. Serve para confirmar que a numeração
de páginas no cabeçalho aparece em crp1, que o cabeçalho usa
a fonte \texttt{\textbackslash Light} em tamanho \texttt{small},
e que a régua separadora do cabeçalho está na cor correta.

\end{document}
